\documentclass[conference]{IEEEtran}
\usepackage{amsmath,amssymb,amsfonts}
\usepackage{graphicx}
\usepackage{url}
\usepackage{booktabs}

\title{Mathematical Derivations for Rusty-Machine: High-Fidelity ML Primitives on CUDA}

\author{\IEEEauthorblockN{Technical Appendix}}

\begin{document}

\maketitle

\section{Introduction}
This document provides the formal mathematical derivations for the machine learning primitives implemented in the \textit{rusty-machine} framework. We focus on the transition from continuous analytical forms to the discrete, vectorized representations executed on NVIDIA CUDA cores using the cuBLAS and cuSOLVER backends.

\section{Linear Regression and Ordinary Least Squares (OLS)}
For a linear model $y = X\theta + \epsilon$, where $X \in \mathbb{R}^{N \times D}$ is the feature matrix and $y \in \mathbb{R}^{N \times 1}$ is the target vector, the objective is to minimize the Residual Sum of Squares (RSS):
\begin{equation}
J(\theta) = \|y - X\theta\|^2 = (y - X\theta)^T(y - X\theta)
\end{equation}

\subsection{Analytical Solution (Normal Equation)}
To find the optimal $\theta$, we set the gradient with respect to $\theta$ to zero:
\begin{align}
\nabla_\theta J(\theta) &= \nabla_\theta (y^Ty - 2\theta^TX^Ty + \theta^TX^TX\theta) \\
&= -2X^Ty + 2X^TX\theta = 0
\end{align}
This yields the Normal Equation:
\begin{equation}
\theta = (X^TX)^{-1}X^Ty
\end{equation}
In \textit{rusty-machine}, this is implemented using \texttt{cuSOLVER}'s Cholesky decomposition ($LL^T$) to solve the linear system for high-dimensional feature spaces.

\section{Logistic Regression and cross-Entropy Loss}
Logistic regression maps the linear output to a probability space using the sigmoid function $\sigma(z)$:
\begin{equation}
h_\theta(x) = \sigma(X\theta) = \frac{1}{1 + e^{-X\theta}}
\end{equation}

\subsection{Cost Function}
The objective is to minimize the Binary Cross Entropy (BCE) loss:
\begin{equation}
J(\theta) = -\frac{1}{N} \sum_{i=1}^N [y_i \log(h_\theta(x_i)) + (1-y_i) \log(1-h_\theta(x_i))]
\end{equation}

\subsection{Gradient Derivation}
The derivative of the sigmoid function is $\sigma'(z) = \sigma(z)(1-\sigma(z))$. Using the chain rule:
\begin{align}
\frac{\partial J}{\partial \theta_j} &= \frac{1}{N} \sum_{i=1}^N (h_\theta(x_i) - y_i)x_{ij} \\
\nabla_\theta J(\theta) &= \frac{1}{N} X^T (\sigma(X\theta) - y)
\end{align}
In \textit{rusty-machine}, the term $X\theta$ is computed via \texttt{cuBLAS SGEMV}, and the element-wise sigmoid and subtraction are handled by custom PTX kernels to maintain zero-copy integrity.

\section{Momentum-Based Optimization}
To accelerate convergence, we implement a velocity-based momentum update. Let $g_t = \nabla_\theta J(\theta_t)$ be the gradient at step $t$:
\begin{align}
v_{t+1} &= \mu v_t + \eta g_t \\
\theta_{t+1} &= \theta_t - v_{t+1}
\end{align}
where $\mu \in [0, 1]$ is the momentum coefficient (typically 0.9) and $\eta$ is the learning rate.

On CUDA architectures, these updates are captured within a \textit{CUDA Graph} to eliminate the Host-Device launch latency of individual vector additions ($v + g$) and subtractions ($\theta - v$).

\section{Memory Lifecycle Mathematics}
The zero-copy mechanism relies on affine pointer mappings. Given a host-resident pointer $P_h$ and its corresponding device-resident allocation $P_d$:
\begin{equation}
P_d = \text{cudaMalloc}(Size(X))
\end{equation}
In Rust, the lifetime $\mathcal{L}$ of $P_d$ is strictly bound such that:
\begin{equation}
\forall t \in [t_{start}, t_{end}], \text{Ownership}(P_d) \in \text{RustBorrowChecker}
\end{equation}
This ensures that the raw memory address $P_d$ passed to the CUDA Stream remains immutable for the duration of the PTX execution cycle.

\end{document}
